\subsection{Two Party $L_2$ Sampling}\label{sec:l2}
In this section we consider the two-party $L_2$ sampling
functionality.
Given input vectors $\vec{w}_1, \vec{w}_2$,
this functionality samples from the distribution
$D_{L_2}(\vec{w}_1, \vec{w}_2)$ with the following
probability mass function:
$$\Pr[ D_{L_2}(\vec{w}_1, \vec{w}_2) \mbox{ samples } i ] 
= \frac{(w_{1,i} + w_{2,i})^2}{ \sum_j (w_{1,j} + w_{2,j})^2 }
= \frac{(w_{1,i} + w_{2,i})^2}{ \| \w_1 + \w_2 \|_2^2 }. $$

We begin by presenting a non-private protocol for two-party $L_2$
sampling with $\tilde{O}(1)$ communication in
Section~\ref{sec:insecure}, the construction is found in Protocol~\ref{fig:l2approx}. We then show how to implement the protocol
securely in Section~\ref{sec:secure}.

\subsubsection{A non-private $L_2$ sampling protocol with $\tilde{O}(1)$
communication} \label{sec:insecure}

We begin by defining and showing how to sample from a helper
distribution $D_{\mathsf{ignore}}$.

\begin{definition}
For input vectors $\vec{w}_1, \vec{w}_2$,
let $D_{\mathsf{ignore}}(\vec{w}_1, \vec{w}_2)$ be the distribution that ``ignores'' the cross term in $D_{L_2}(\vec{w}_1, \vec{w}_2)$.
I.e.~$D_{\mathsf{ignore}}(\vec{w}_1, \vec{w}_2)$
samples index $i \in [n]$
with probability $\frac{w^2_{1,i} + w^2_{2,i}}{||\vec{w}_1||_2^2 + ||\vec{w}_2||_2^2}$.
\end{definition}

\begin{lemma}\label{lem:ignore}
There exists a protocol $\Pi_{\mathsf{ignore}}$ for sampling from
$D_{\mathsf{ignore}}(\vec{w}_1, \vec{w}_2)$ with $\tilde{O}(1)$
communication.
\end{lemma}

\begin{proof}
Let $\vec{w'}_b = (w_{b,1}^2, \ldots, w_{b,n}^2).$
The lemma follows by observing the following:
$$D_\ig(\w_1, \w_2) = D_{L_1}(\vec{w'}_1, \vec{w'}_2).$$ \qed
%Each party samples from its own distribution (i.e., $P_b$ outputs $i$
%with probability $\frac{w_{b,i}^2}{||\vec{w}_b||^2}$).  Then the
%parties do a weighted coin flip to choose which of these two values is
%output ($P_b$'s value is chosen with probability
%$\frac{||\vec{w}_b||_2^2}{||\vec{w}_1||_2^2+||\vec{w}_2||_2^2}$.)  It
%is easy to see that this protocol only requires $\tilde{O}(1)$
%communication\anote{Should we call this logarithmic?} and produces the
%correct output distribution.
\end{proof}

\begin{definition}
For $i \in [n]$, let the corrective parameter function be defined as 
$$f_c(\vec{w}_1, \vec{w}_2, i) := \frac{w_{1,i}^2 + 2w_{1,i}w_{2,i} +
w_{2,i}^2}{||\vec{w}_1||_2^2 + ||\vec{w}_2||_2^2}.$$
\end{definition}


\begin{definition} 
The constant 
$c := c(\vec{w}_1, \vec{w}_2)$ 
is defined as
\[
c(\vec{w}_1, \vec{w}_2) := \frac{||\vec{w}_1 + \vec{w}_2||_2^2}{||\vec{w}_1||_2^2 + ||\vec{w}_2||_2^2}
\]
This ensures that 
for every $i$,
$f_c(\vec{w}_1, \vec{w}_2, i) = c \cdot \Pr_{D_{L_2}(\vec{w}_1, \vec{w}_2)}[i]$.
\end{definition}

\iffalse
\begin{proof}
Define 
$c(\vec{w}_1, \vec{w}_2) := \frac{||\vec{w}_1 + \vec{w}_2||_2^2}{||\vec{w}_1||_2^2 + ||\vec{w}_2||_2^2}$.
Clearly, by definition of $f_c$ and $D_{L_2}(\vec{w}_1, \vec{w}_2)$ it is the case that
$f_c(\vec{w}_1, \vec{w}_2, i) = c(\vec{w}_1, \vec{w}_2) \cdot \Pr_{D_{L_2}(\vec{w}_1, \vec{w}_2)}[i]$.
It remains to show that $c(\vec{w}_1, \vec{w}_2) \in [1,2]$.
For the upper bound,
\begin{align*}
    c(\vec{w}_1, \vec{w}_2) &= \frac{||\vec{w}_1||^2_2 + 2\langle \vec{w}_{1}, \vec{w}_{2} \rangle + ||\vec{w}_2||_2^2}{||\vec{w}_1||_2^2 + ||\vec{w}_2||_2^2}\\
    &\leq \frac{2(||\vec{w}_1||^2_2 + ||\vec{w}_2||_2^2)}{||\vec{w}_1||_2^2 + ||\vec{w}_2||_2^2}\\
    &= 2.
\end{align*}
For the lower bound,  $c(\vec{w}_1, \vec{w}_2) \geq 1$, since $2\langle \vec{w}_1, \vec{w}_2 \rangle \geq 0$ (since in our setting $\vec{w}_1, \vec{w}_2$ are positive).
\end{proof}
\fi

The following lemma will be useful for arguing the validity of the
final protocol.

\begin{lemma} \label{lem:close_dist}
For all $i \in \mathsf{supp}(D_{L_2}(\vec{w}_1, \vec{w}_2))$, 
$\Pr_{D_{L_2}(\vec{w}_1, \vec{w}_2)}[i] \leq 2/c \cdot \Pr_{D_{\mathsf{ignore}}(\vec{w}_1, \vec{w}_2)}[i]$.
\end{lemma}

\begin{proof}
\begin{align*}
    \Pr_{D_{L_2}(\vec{w}_1, \vec{w}_2)}[i] &= \frac{w^2_{1,i} + 2w_{1,i}w_{2,i} + w^2_{2,i}}{||\vec{w}_1||_2^2 + 2\langle \vec{w}_1, \vec{w}_2 \rangle + ||\vec{w}_2||_2^2}\\
    &=  \frac{w^2_{1,i} + 2w_{1,i}w_{2,i} + w^2_{2,i}}{c \cdot (||\vec{w}_1||_2^2  + ||\vec{w}_2||_2^2)}\\
    &\le \frac{2 \cdot (w_{1,i}^2  + w_{2,i}^2)}{c \cdot ||\vec{w}_1||_2^2  + ||\vec{w}_2||_2^2}\\
    & = \frac{2}{c} \cdot \Pr_{D_{\mathsf{ignore}}(\vec{w}_1, \vec{w}_2)}[i]
\end{align*}
The inequality holds since
\begin{align*}
2(w^2_{1,i} + w^2_{2,i}) - (w^2_{1,i} + 2w_{1,i}w_{2,i} + w^2_{2,i}) 
&= w^2_{1,i} -2w_{1,i}w_{2,i} + w^2_{2,i}\\
&= (w_{1,i} - w_{2,i})^2\\
&\geq 0.
\end{align*}
\qed
\end{proof}

We now present the $L_2$ sampling protocol $\Pi_{L_2}$, which is described in Protocol~\ref{fig:l2approx}. We show the correctness and efficiency of the protocol. 

\floatname{algorithm}{Protocol}
\begin{algorithm}[htbp]

\paragraph{Inputs:} Parties $P_1$ and $P_2$ have inputs $\w_1$ and $\w_2$
respectively. 

\medskip

The protocol proceeds as follows:
\BEN
    \item Parties run $\Pi_{\mathsf{ignore}}$ 
    with inputs $\vec{w}_1, \vec{w}_2$
    that samples from $D_{\mathsf{ignore}}(\vec{w}_1, \vec{w}_2)$ and obtain output $i$.
    
    \item For $b \in \{1, 2\}$, $P_b$ sends $w_{b,i}, ||\vec{w}_b||_2^2$.
    Both parties compute 
    \[
    \Pr_{D_{\mathsf{ignore}}(\vec{w}_1, \vec{w}_2)}[i] = \frac{w_{1,i}^2 + w_{2,i}^2}{||\vec{w}_1||_2^2 + ||\vec{w}_2||_2^2}
    \quad
    \mbox{and}
    \quad
    f_c(\vec{w}_1, \vec{w}_2, i) = \frac{w_{1,i}^2 + 2w_{1,i}w_{2,i} + w_{2,i}^2}{||\vec{w}_1||_2^2 + ||\vec{w}_2||_2^2}
    \]
    \item Parties output $i$ with probability 
    \begin{align*}
        \frac{f_c(\vec{w}_1, \vec{w}_2, i)}{2 \cdot \Pr_{D_{\mathsf{ignore}}(\vec{w}_1, \vec{w}_2)}[i]} &= \frac{c \cdot \Pr_{D_{L_2}(\vec{w}_1, \vec{w}_2)}[i]}{2 \cdot \Pr_{D_{\mathsf{ignore}}(\vec{w}_1, \vec{w}_2)}[i]}\\
        &= \frac{\Pr_{D_{L_2}(\vec{w}_1, \vec{w}_2)}[i]}{2/c \cdot \Pr_{D_{\mathsf{ignore}}(\vec{w}_1, \vec{w}_2)}[i]}
        \end{align*}
    and otherwise return to step 1.
    
\EEN
\caption{Protocol for exact $L_2$ sampling ($\Pi_{L_2}$)}
\label{fig:l2approx}
\end{algorithm}

\begin{lemma} \label{lem:insecure_L2_prot}
With all but negligible probability, on inputs $\vec{w}_1, \vec{w}_2$, $\Pi_{L_2}$ samples exactly correctly from $D_{L_2}(\vec{w}_1, \vec{w}_2)$, and has communication
$\tilde{O}(1)$.
\end{lemma}

\begin{proof}
Note that $\Pi_{L_2}$ simply performs rejection sampling in a distributed setting where sampling from $D_{\mathsf{ignore}}(\vec{w}_1, \vec{w}_2)$ and computing the probabilities is done in a distributed manner. It is therefore well-known that
as long as for all $i \in [n]$,
\begin{equation} \label{eq:rejection}
    \Pr_{D_{L_2}(\vec{w}_1, \vec{w}_2)}[i] \leq 2/c \cdot \Pr_{D_{\mathsf{ignore}}(\vec{w}_1, \vec{w}_2)}[i],
\end{equation}
then $\Pi_{L_2}$ samples from the exact correct distribution, and
the number of samples required from $D_{\mathsf{ignore}}(\vec{w}_1, \vec{w}_2)$ in protocol $\Pi_{L_2}$ follows a geometric distribution with probability $c/2$.
Thus, if condition (\ref{eq:rejection}) is met, the protocol 
samples exactly correctly and completes in an expected $2/c$ (with $2/c \leq 2$, since $c \ge 1$) number of rounds.   
Further, it can be immediately noted that condition (\ref{eq:rejection}) is met due to Lemma~\ref{lem:close_dist}.
Finally, each round has $\tilde{O}(1)$ communication, since $\Pi_{\mathsf{ignore}}$ has communication $\tilde{O}(1)$ (by Lemma~\ref{lem:ignore})
and since, in addition to that, only a constant number of length $\tilde{O}(1)$ values are exchanged in each round.
Combining the above, we have that $\Pi_{L_2}$ has expected
communication $\tilde{O}(1)$
and worst case (with all but negligible probability) communication $\tilde{O}(1)$.
\qed

\end{proof}

\begin{remark}
Note that the protocol and analysis above
did not require that vectors $\vec{w}_1, \vec{w}_2$ are normalized.
I.e.~we do not require that $||\vec{w}_1||_1$
or $||\vec{w}_2||_1$ are equal to $1$
or to each other.
\end{remark}

\subsubsection{Secure $L_2$ Sampling From FHE}
\label{sec:secure}

\floatname{algorithm}{Protocol}
\begin{algorithm}[h]
\textbf{Inputs:} Party $P_b$ has input $\vec{w}_b$. 

\BEN
\item Let $B \in \tilde{O}(1)$. The parties perform
the following steps for $j \in [B]$:

\begin{enumerate}
\item Sample from $D_{\mathsf{ignore}}(\vec{w}_1,\vec{w}_2)$
by doing the following:
Invoke ideal functionality $\F_{L_1}^{ss}$ with $P_1$'s input set to $\vec{w}_1 \odot \vec{w}_1$
and $P_2$'s input 
set to $\vec{w}_2 \odot \vec{w}_2$. 
Let $\langle i_j \rangle$ be the secret share of the output
index. 

\item Parties compute encryptions of $w_{1,i_j}, w_{2,i_j}$ using a threshold FHE
scheme as follows.  
\begin{itemize}\item Parties compute an encryption of $i_j$ by exchanging encryptions of their shares and adding them.  \item Party $b$ encrypts $\vec{w}_b$ and uses FHE to locally compute an encryption of $w_{b,i_j}$. \item  The parties then send these ciphertexts to each other.\end{itemize}
\anote{Commenting out this footnote unless we can figure out how to include it in a protocol box}
% \footnote{A simpler version of this step can be realized by having $\F_{L_1}$ output a secret sharing of $w_{b,i_j}$ in addition to $i_j$, but we present it this way for modularity purposes.}

\item Rejection Sampling. Compute a threshold FHE ciphertext $\widehat{\mathsf{bias}}_j$ that
encrypts
\[
\frac{f_c(\vec{w}_1, \vec{w}_2)_{i_j}}{2 \cdot \Pr_{D_{\mathsf{ignore}}(\vec{w}_1, \vec{w}_2)}[i_j]} =
\frac{w_{1,i_j}^2 + 2w_{1,i_j}w_{2,i_j}+w_{2,i_j}^2}{2(w_{1,i_j}^2 + w_{2,i_j}^2)}.
\]
Invoke ideal functionality $\F_{2PC}$ 
that takes  encrypted bias $\widehat{\mathsf{bias}}_j$, the threshold decryption keys, index $i_j$, and random bits. The functionality executes a circuit that flips a coin with bias $\widehat{\mathsf{bias}}_j$ and returns a ciphertext $\widehat{\mathsf{out}}_j$, which is
an encryption of $i_j$ if the coin evaluates to $1$
and an encryption of $0$ otherwise.
\end{enumerate}

\item Execute $\F_{2PC}$ for the following circuit: 
      \BEN
      \item Input: $(\widehat{\mathsf{out}}_1, \ldots, \widehat{\mathsf{out}}_B)$ and threshold decryption keys. 
      \item Output: $i_j$ corresponding to the minimum
      $j$ such that $\widehat{\mathsf{out}}_j$ decrypts to $i_j \neq 0$.
      Or $\bot$ if no such $j \in [B]$ exists.
      \EEN
\EEN

\caption{Two-party $L_2$ sampling in the $(\F_{L_1}, F_{2PC})$-hybrid.}
\label{fig:l2}
\end{algorithm}

\paragraph{$L_2$ sampling protocol.}
We present our secure $L_2$ sampling protocol in
Protocol~\ref{fig:l2}.
For two $n$-dimensional vectors $\vec{w}_1$ and $\vec{w}_2$, we
denote by $\vec{w}_1 \odot \vec{w}_2$ the $n$-dimensional
vector whose $i$-th entry is equal to
$w_{1,i} \cdot w_{2,i}$.

Our $L_2$ sampling protocol uses ideal functionality $\F_{L_1}^{ss}$, which works essentially the same as $\F_{L_1}$ except that the output index is secret shared among both parties. We can securely realize this functionality with semi-honest security through a trivial change in the protocol $\Pi_{L_1}$;~for the sake of completeness, we provide the details in Appendix~\ref{apx:ssL1}.
 

\paragraph{Efficiency and correctness.}
It is clear that the total communication complexity of the protocol is $\tilde{O}(1)$, since each step in the loop has complexity $\tilde{O}(1)$ and the loop iterates $B \in \tilde{O}(1)$ number of times.
Correctness is also immediate, since the protocol simply implements the $\Pi_{L_2}$ sampling procedure, which was proven in Section \ref{sec:insecure} to be correct, and to require at most $B \in \tilde{O}(1)$ samples, with all but negligible probability,
 
 \paragraph{Security.} Security of our protocol is stated through the following theorem.
\begin{theorem}
\label{th:l2}
Assuming the existence of threshold FHE with
$\mathsf{IND}\mbox{-}\mathsf{CPA}$ security,
Protocol~\ref{fig:l2} securely realizes the $L_2$ sampling
functionality in
the $\{\F_{L_1}^{ss}, \F_{2PC}\}$-hybrid model with semi-honest security.
\end{theorem}
%
We provide the proof in Appendix~\ref{apx:pfTheorem3}.
